\section{Limitations}\label{sec:limitations}
The feasibility result has a deliberately narrow statistical scope. Proposition~3 applies to deterministic, uniformly valid, distribution-free upper confidence bounds based only on iid bounded category losses. It does not rule out nonvacuous certification from fewer categories under pre-specified parametric or hierarchical models, informative side data, or randomized procedures. Such approaches exchange distribution-free generality for additional assumptions and require their own coverage analyses. The counts 14, 29, and 59 are optimistic necessary limits without multiplicity, not sufficient sample sizes for CRESS. The frozen family allocation and declared Hoeffding rule require more.

The category interpretation also depends on the unit definition. The strict audit contains only three or four certification categories, and repeated images, corruption views, or support seeds do not constitute additional independent category draws. Image-unit certificates under the selected-mixture iid model therefore address a different estimand rather than a partial certificate for a new category. CRESS certifies source-domain risk under the assumptions of Proposition~1. The benchmark categories are a fixed heterogeneous collection, so their interpretation as independent draws from a source-category meta-population is a modeling assumption rather than an empirically testable fact in these data. Applying the selected threshold to a particular target additionally requires source-risk dominance. Category exchangeability supports only risk marginalized over a random target-category draw; it does not control every realized category. Support-only normalization and source-view selection establish neither condition. Matched source selection also assumes reliable condition metadata.

The empirical evidence is limited to MVTec, VisA, and MPDD. The controlled shift study covers Gaussian noise, blur, brightness/contrast, and JPEG corruption, not the full range of temporal drift, sensor replacement, process changes, correlated failures, or deployment prevalence. The empirical FAR of 0.341 is the largest dataset-level aggregate in this frozen synthetic-corruption grid, not a population failure probability or a universal response to Gaussian noise. The image-unit sensitivity additionally assumes independent source-image draws; repeated or correlated production images would reduce its effective evidence. Prospective factory evidence is needed to determine whether the same failure mode and category-count assumptions hold under real operational shifts.

Finally, the frozen DINOv2 PCA residual is an inherited ranking substrate. Official AnomalyDINO and reported WinCLIP results follow different protocols and serve as external references; the paper does not claim ranking state of the art. LOIO calibration and full-support test residuals remain asymmetric, and the pooled discrete-grid audit diagnoses rather than restores exchangeability. Reported scoring latency uses cached DINOv2 features and excludes backbone inference, preprocessing, and input/output overhead.
